\chapter{Blood in Semen}
\index{Blood in Semen}

\section*{Definition and Overview}
Blood in semen (haematospermia) refers to the presence of blood in the ejaculate, which appears pink, red, or brown. The blood almost always originates from bleeding somewhere in the male reproductive tract -- the prostate gland, seminal vesicles, epididymis, or urethra. In younger men the cause is overwhelmingly benign, most often an infection or minor inflammation, and the bleeding typically clears spontaneously. However, because it can occasionally signal a more serious underlying problem, particularly in men over 40, blood in semen should always be assessed by a doctor rather than ignored.

\section*{Epidemiology}
Haematospermia is a fairly common reason for men to seek medical advice, although the true frequency is hard to measure because many men never mention the symptom or the episode is a single, self-limited event. In men under 40 the great majority of cases are benign and resolve without treatment. In older men the risk of a significant underlying cause rises, and the symptom can be the first pointer to prostate disease, including prostate cancer. Most episodes are single or brief, but recurrent or persistent haematospermia warrants structured investigation.

\section*{Aetiology and Pathophysiology}
Blood in semen arises when bleeding occurs at any point along the pathway that produces and transports semen:

\begin{itemize}
    \item \textbf{Infections:} prostatitis, seminal vesiculitis, urethritis, or epididymitis; sexually transmitted infections such as chlamydia and gonorrhoea can cause urethral bleeding
    \item \textbf{Prostate problems:} benign prostatic enlargement, prostatic inflammation, or prostate stones
    \item \textbf{Procedures:} prostate biopsy, cystoscopy, or other urinary tract instrumentation
    \item \textbf{Trauma:} vigorous sexual activity, masturbation, or injury to the perineum
    \item \textbf{Prolonged abstinence:} accumulation of secretions that are released with the next ejaculation
    \item \textbf{Blood-thinning medicines:} anticoagulants and antiplatelet agents make bleeding more likely
    \item \textbf{Rarely, malignancy:} prostate, seminal vesicle, or testicular cancer, particularly in older men
\end{itemize}

\section*{Clinical Features}
\begin{itemize}
    \item Pink, red, or brown colouring or streaks in the ejaculate; a brown tint suggests older blood
    \item Often a single episode with no other symptoms
    \item Pain or burning during ejaculation or urination (suggests infection)
    \item Pelvic, perineal, or testicular discomfort (prostatitis or epididymitis)
    \item Urinary symptoms: urgency, frequency, weak stream, or hesitancy (prostate problems)
    \item Fever and chills with acute infection
    \item A palpable testicular lump or scrotal swelling
\end{itemize}

\section*{Differential Diagnosis}
\begin{itemize}
    \item Prostatitis or seminal vesiculitis
    \item Sexually transmitted infections, including chlamydia and gonorrhoea
    \item Benign prostatic enlargement or prostate stones
    \item Epididymitis or orchitis
    \item Recent prostate biopsy or other urological procedure
    \item Prostate, bladder, or testicular cancer (less common)
    \item Bleeding from urinary rather than reproductive sources, such as blood in the urine mistaken for haematospermia
\end{itemize}

\section*{When to Seek Medical Care}
See a doctor for any episode of blood in semen, and especially if it is recurrent or persistent, occurs after age 40, or is accompanied by pain, fever, urinary symptoms, or a scrotal lump. Men taking anticoagulants, and those with risk factors for prostate cancer, should be assessed promptly.

\textbf{Contact your local emergency services for:}
\begin{itemize}
    \item Inability to pass urine
    \item Very heavy bleeding, or blood in the urine with large clots
    \item High fever with chills and severe pelvic or testicular pain
    \item Sudden, severe pain or swelling of the testicle (possible torsion)
\end{itemize}

\section*{Diagnosis and Investigations}
\begin{itemize}
    \item \textbf{History:} duration and frequency of bleeding, associated pain, urinary symptoms, recent procedures, sexual history, and current medicines
    \item \textbf{Physical examination:} examination of the abdomen, testes, and a digital rectal examination of the prostate
    \item \textbf{Urine tests:} urinalysis and urine culture to detect infection or blood
    \item \textbf{STI screening:} swab or urine testing for sexually transmitted infections
    \item \textbf{Blood tests:} prostate-specific antigen (PSA) in older men; inflammatory markers if fever is present
    \item \textbf{Imaging:} scrotal or transrectal ultrasound in selected cases
    \item \textbf{Referral:} to a urologist for cystoscopy or cross-sectional imaging if symptoms persist or a serious cause is suspected
\end{itemize}

\section*{Management}
\begin{itemize}
    \item \textbf{Reassurance:} in young men with a single episode and no other symptoms, no specific treatment is usually needed
    \item \textbf{Infection:} antibiotics for prostatitis or sexually transmitted infections, with partner treatment where required
    \item \textbf{Prostate problems:} treatment of benign enlargement under medical advice
    \item \textbf{Medicines:} discuss blood-thinning medicines with a doctor before stopping or adjusting them
    \item \textbf{Activity:} reduce vigorous sexual activity or perineal trauma while bleeding settles
    \item \textbf{Urological care:} further investigation and treatment if bleeding persists or a serious cause is found
    \item \textbf{Follow-up:} a repeat review is sensible to confirm the bleeding has stopped
\end{itemize}

\section*{Complications}
\begin{itemize}
    \item Anxiety and distress about what the symptom might mean
    \item Untreated infection spreading within the reproductive or urinary tract
    \item A missed serious underlying condition, particularly prostate cancer in older men
    \item Iron deficiency anaemia is rare but possible with very frequent or heavy bleeding
\end{itemize}

\section*{Prognosis and Outcomes}
The outlook is generally excellent. In most men, especially those under 40, blood in semen resolves on its own within a few weeks and never recurs. When an infection is identified and treated, the bleeding usually stops quickly. Serious causes are uncommon, but when present they are best treated early, which is why medical assessment matters. Men over 40 and those with recurrent bleeding are most likely to benefit from urological investigation.

\section*{Prevention and Self-Care}
\begin{itemize}
    \item Practise safe sex to reduce the risk of sexually transmitted infections
    \item Seek early treatment for urinary symptoms or pelvic discomfort
    \item Avoid trauma to the genital area during vigorous activity
    \item Keep blood-thinning medicines under medical review and report bleeding episodes to the prescribing doctor
    \item Attend recommended prostate checks if you are over 40 or at increased risk
    \item Do not assume the symptom will resolve on its own -- have it assessed
\end{itemize}

\section*{Sources and Further Reading}
The following sources provide reliable, up-to-date information on blood in semen and male reproductive health:

\begin{itemize}
    \item NHS. \textit{NHS guidance on blood in semen and urological symptoms.} https://www.nhs.uk/conditions/
    \item Mayo Clinic. \textit{Information on haematospermia and prostate health.} https://www.mayoclinic.org/
    \item MedlinePlus. \textit{Blood in semen patient information.} https://medlineplus.gov/
    \item MSD Manual. \textit{Overview of seminal vesicle and prostate disorders.} https://www.msdmanuals.com/
    \item Centers for Disease Control and Prevention. \textit{Sexually transmitted infections information.} https://www.cdc.gov/
    \item NICE. \textit{Suspected cancer recognition and referral guidelines.} https://www.nice.org.uk/
\end{itemize}
